% Figure: a snapshot of a linearly polarised plane wave travelling in +z.
% Plotted are the in-phase electric field E (arbitrary units), the magnetic
% field H = E/eta (same phase, smaller amplitude in these units), and the
% instantaneous Poynting magnitude S = E*H, which is everywhere non-negative
% and oscillates at twice the field frequency. The point: E and H are in
% phase for a travelling wave, so power flows in pulses, and the time
% average of S is half its peak.
\begin{figure}[htbp]
\centering
\begin{tikzpicture}
\begin{axis}[
  width=12cm, height=6cm,
  xlabel={position $z$ (wavelengths)},
  ylabel={amplitude (arb.\ units)},
  xmin=0, xmax=2, ymin=-1.2, ymax=1.2,
  domain=0:2, samples=300,
  axis lines=middle,
  enlargelimits=false,
  xtick={0,0.5,1,1.5,2},
  ytick={-1,-0.5,0,0.5,1},
  legend pos=south east,
  legend cell align=left,
]
  % E field, unit amplitude
  \addplot[engblue, very thick] { sin(deg(2*pi*x)) };
  \addlegendentry{$E$}
  % H field, in phase, 0.6 amplitude for visibility (E/eta scaled)
  \addplot[engred, very thick, dashed] { 0.6*sin(deg(2*pi*x)) };
  \addlegendentry{$H = E/\eta$}
  % instantaneous Poynting magnitude S = E*H, always >= 0, scaled
  \addplot[engorange, thick] { 0.6*(sin(deg(2*pi*x)))^2 };
  \addlegendentry{$S = EH$}
  % time-average line of S = 0.3 (half of 0.6 peak)
  \draw[enggray, dashed] (axis cs:0,0.3) -- (axis cs:2,0.3);
  \node[enggray, font=\scriptsize, anchor=south west] at (axis cs:0.02,0.30) {$\langle S\rangle$};
\end{axis}
\end{tikzpicture}
\caption[Plane wave fields and Poynting flux in phase]{Snapshot of a linearly
polarised travelling plane wave. The electric field $E$ and magnetic field
$H = E/\eta$ are in phase and in spatial step, the signature of a
travelling wave as opposed to a standing one. The instantaneous Poynting
magnitude $S = EH$ is the product of two in-phase sinusoids, so it never
goes negative and oscillates at twice the field frequency between zero and a
peak; energy flows in the $+z$ direction in pulses. The time average
$\langle S\rangle$, the dashed line, is half the peak, which is the origin
of the factor $\tfrac{1}{2}$ in $\langle S\rangle = E_0^2/(2\eta)$.}
\label{fig:vol04ch10:poynting-energy}
\end{figure}
